\begin{figure}[htbp]
  \centering
  \textbf{\small (a) Da origem ao facto}\\[3pt]
  \begin{tikzpicture}[x=1cm, y=1cm, font=\scriptsize]
    \node[ntcaixa, text width=2.5cm, minimum height=11mm] (o1) at (1.35,1.45)
      {Historiador de processo\\[1pt]{\tiny \emph{tags}, valor horário}};
    \node[ntcaixa, text width=2.5cm, minimum height=11mm] (o2) at (1.35,0.00)
      {Balanços de utilidades\\e de produção\\[1pt]{\tiny valor diário}};
    \node[ntcaixa, text width=2.5cm, minimum height=11mm] (o3) at (1.35,-1.45)
      {Relatórios e ficheiros\\[1pt]{\tiny valor mensal}};
    \node[ntcaixa, text width=2.9cm, minimum height=15mm] (ag) at (5.4,0.00)
      {Agregação e conversão\\[1pt]
       {\tiny soma horária $\to$ dia; parâmetros lidos de tabelas
        (PCI, entalpias, factores)}};
    \node[ntdestaque, text width=3.3cm, minimum height=15mm] (pr) at (9.5,0.00)
      {Precedência\\linha a linha\\[1pt]
       {\tiny mapeamento mestre: 286 linhas\\
        unidade\,$\times$\,vetor\,$\times$\,papel}};
    \node[ntdestaque, text width=2.4cm, minimum height=15mm] (fa) at (13.2,0.00)
      {Tabelas de factos\\[1pt]{\tiny com origem, \emph{flags} e versão}};
    \draw[ntseta] (o1.east) -- ++(0.35,0) |- (ag.west);
    \draw[ntseta] (o2.east) -- (ag.west);
    \draw[ntseta] (o3.east) -- ++(0.35,0) |- (ag.west);
    \draw[ntsetaf] (ag) -- (pr);
    \draw[ntsetaf] (pr) -- (fa);
    \node[ntmini, anchor=north, align=center] at (9.5,-1.05)
      {nenhuma constante de conversão\\vive fora de uma tabela versionada};
  \end{tikzpicture}

  \vspace{4mm}
  \textbf{\small (b) Esquema em estrela simplificado}\\[3pt]
  \begin{tikzpicture}[x=1cm, y=1cm, font=\scriptsize]
    \node[ntcaixa, text width=2.1cm, minimum height=9mm] (dc) at (1.20,1.65)
      {\textsc{dim} Calendário\\[1pt]{\tiny dia}};
    \node[ntcaixa, text width=2.1cm, minimum height=9mm] (du) at (1.20,0.15)
      {\textsc{dim} Unidade\\[1pt]{\tiny 16 unidades}};
    \node[ntcaixa, text width=2.1cm, minimum height=9mm] (dv) at (1.20,-1.35)
      {\textsc{dim} Vetor\\[1pt]{\tiny 7 vetores}};

    \node[ntdestaque, text width=4.3cm, minimum height=12mm] (fe) at (6.45,1.65)
      {\textsc{fact} Energia\\[1pt]
       {\tiny dia\,$\times$\,unidade\,$\times$\,vetor --- 200\,924 linhas,
        2020-01-01 a 2026-07-31}};
    \node[ntdestaque, text width=4.3cm, minimum height=12mm] (fc) at (6.45,0.15)
      {\textsc{fact} Carga\\[1pt]
       {\tiny dia\,$\times$\,unidade --- 45\,408 linhas, total e componentes}};
    \node[ntcaixa, text width=4.3cm, minimum height=12mm] (fp) at (6.45,-1.35)
      {\textsc{fact} Propriedades da carga\\[1pt]
       {\tiny dia\,$\times$\,unidade --- 19\,105 linhas, 8 unidades}};

    \foreach \f in {fe,fc,fp} {
      \draw[corCinzaC, line width=.5pt] (dc) -- (\f);
      \draw[corCinzaC, line width=.5pt] (du) -- (\f);
    }
    \draw[corCinzaC, line width=.5pt] (dv) -- (fe);

    \node[ntcaixa, text width=3.0cm, minimum height=14mm, fill=corFundo] (el)
      at (11.75,1.65)
      {\textbf{Eletricidade}\\[1pt]
       {\tiny medida ao mês e alocada ao dia por regra declarada: a variação
        diária que se lê é a do repartidor, não a do consumo}};
    \draw[ntsetat, draw=corV3] (el.west) -- (fe.east);

    \node[ntcaixa, text width=3.0cm, minimum height=10mm] (pa) at (11.75,-0.55)
      {Tabelas de parâmetros\\[1pt]{\tiny PCI, entalpias, limiares de carga}};
    \draw[corCinzaC, line width=.5pt] (pa.west) -- (fc.east);
  \end{tikzpicture}
  \caption[Arquitetura de dados: proveniência e granularidade]%
  {Arquitetura de dados da plataforma. No painel~(a), o percurso de cada
   valor desde a origem até ao facto: a agregação e a conversão usam
   parâmetros lidos de tabelas, e a escolha entre canais é resolvida linha
   a linha pelo mapeamento mestre, não por código. No painel~(b), o
   esquema em estrela reduzido às três tabelas de factos e às três
   dimensões que este trabalho usa, com a granularidade de cada facto
   escrita por extenso. A eletricidade é a única exceção à granularidade
   diária: é medida ao mês e alocada ao dia por uma regra declarada, pelo
   que a variação diária que dela se lê não é variação de consumo. O
   modelo completo tem 32 tabelas e está descrito no dicionário de dados
   (Apêndice~\ref{app:manual}).}
  \label{fig:arquitetura-dados}
\end{figure}
